\section{Phase 4: on a real chip (the Zybo FPGA board)}\label{sec:fpga}

Everything so far ran on computers \emph{pretending} to be the Ibex: Verilator simulates the chip's circuit
tick by tick. A simulation is only as good as its model, so the last step is to build the real thing and
check that it behaves the same.

\begin{kid}
A flight simulator is great for learning, but one day you have to fly a real plane. If the real plane does
exactly what the simulator said, you can trust everything you learned in the simulator. That is what this
phase checks, and it also finds out what the simulator could not tell us.
\end{kid}

\subsection{What an FPGA is}

An \textbf{FPGA} (field-programmable gate array) is a chip full of small, unconnected building blocks: tiny
look-up tables (\textbf{LUTs}) that can act as any small logic gate, flip-flops that remember one bit, blocks of
memory (\textbf{block RAM}) and wires that can be switched on and off. A file called a \textbf{bitstream} tells
the chip how to connect them. Load a different bitstream and the same chip becomes a different circuit. So we can
load ``Ibex with our LOD instruction'' into it, and it then \emph{is} that processor, running at a real clock.

We used the \textbf{Zybo} board from Digilent. Its chip, a Xilinx \textbf{Zynq XC7Z010}, has two parts:
\begin{itemize}
  \item a normal ARM processor (the \emph{processing system}), which runs a small control program, talks to the
        PC and owns the 512~MB DDR3 memory of the board;
  \item the programmable logic (the FPGA part), with 17{,}600 LUTs, 35{,}200 flip-flops and 60 block RAMs of
        4~KB each (240~KB in total). This is where our Ibex lives.
\end{itemize}

The tools are Xilinx \textbf{Vivado} (turns the hardware description into a bitstream and reports area, timing
and power) and \textbf{Vitis} (compiles the ARM program).

\subsection{The system we built}

\begin{figure}[H]
\centering
\begin{tikzpicture}[
  bx/.style={draw,rounded corners,minimum height=1.1cm,text width=2.9cm,align=center,font=\scriptsize},
  lb/.style={font=\tiny,fill=white,inner sep=1pt},
  >={Stealth[length=1.6mm]}]
  \node[bx,fill=gray!10] (pc) at (-0.7,0) {\textbf{PC}\\\file{host/run\_board.py}};
  \node[bx,fill=blue!6] (arm) at (4.6,0) {\textbf{ARM Cortex-A9}\\loads programs,\\collects results};
  \node[bx,fill=orange!10] (ibex) at (9.9,0) {\textbf{Ibex RV32IC + LOD}\\no multiplier, 50~MHz};
  \node[bx,fill=blue!3] (ddr) at (4.6,-2.4) {\textbf{DDR3}, 512~MB\\about 15 cycles per access};
  \node[bx,fill=orange!5] (bram) at (9.9,-2.4) {\textbf{block RAM}, 192~KB\\1 cycle, like the simulator};
  \draw[<->] (pc) -- node[lb,above=2pt]{USB serial} (arm);
  \draw[<->] (arm) -- node[lb,above=2pt]{control, loading} (ibex);
  \draw[<->] (arm) -- node[lb,left=2pt]{its own memory} (ddr);
  \draw[<->] (ibex) -- node[lb,right=2pt]{code (+ data)} (bram);
  \draw[<->] (ibex.south west) -- node[lb,sloped,above=1pt]{data in DDR mode} (ddr.north east);
  \draw[dashed,gray] (2.85,0.95) rectangle (11.7,-3.25);
  \node[font=\tiny,gray,anchor=north east] at (11.7,-3.3) {the Zynq chip and its DDR};
\end{tikzpicture}
\caption{The board system (\file{fpga\_zybo/}). Orange: the programmable logic we designed
(\file{rtl/ibex\_zybo\_soc.sv}); blue: the Zynq's ARM side.}
\label{fig:board}
\end{figure}

The Ibex is configured exactly as in the simulator: RV32IC, \textbf{no multiplier}, the LOD instruction added,
and the same memory map, so the very same compiled programs run unchanged. One test run goes like this:

\begin{enumerate}
  \item The PC sends a program (weights, code and 5 test inputs) to the ARM over the USB cable.
  \item The ARM copies it into the Ibex's memory while holding the Ibex in reset, then lets it run.
  \item The Ibex computes the 5 inputs and prints one line per input: prediction, a fingerprint (hash) of all
        output values, and its cycle counter for every layer.
  \item The ARM passes the lines back to the PC, which compares them with the expected answers.
\end{enumerate}

\subsection{The memory problem, and two ways to run}

The simulator gave the Ibex 1~MB of perfect one-cycle memory. The Zybo's block RAM is only 240~KB, and our
largest programs need more than twice that (CIFAR-10 int16 about 600~KB). Two things helped:

\begin{itemize}
  \item \textbf{Sharing memory between layers.} A layer's input is no longer needed once the next layers have
        used it, so buffers whose lifetimes do not overlap can share the same bytes. This shrank the programs a
        lot, and every Speech Commands program now fits in 192~KB.
  \item \textbf{Two memory settings.}
        \textbf{bram}: the whole program in block RAM. This is exactly the simulator's situation, so the cycle
        counts \emph{must} be the same.
        \textbf{ddr}: code in block RAM, all data (weights, activations, tables) in the board's DDR. Every
        program runs this way. It is what a real microcontroller with external memory would see.
\end{itemize}

\begin{kid}
Block RAM is a book on your desk: you read a word instantly. DDR is a book in the basement: every single word
means a trip down the stairs. The work you do with the words is the same; you just spend a lot of time walking.
\end{kid}

\begin{careful}
Sharing buffers changes where the compiler puts some code, and on this CPU that alone moves the cycle count of
a layer by up to a few per cent (Section~\ref{sec:tiny} used the earlier build). So the board's block-RAM runs
are compared with a Verilator run of \emph{exactly the same} program. The speed-ups barely move: Speech
Commands int8 table 1.97$\times$ $\to$ 1.95$\times$, int16 Dn1 2.94$\times$ $\to$ 2.93$\times$.
\end{careful}

\subsection{Did the board agree with the simulation?}

\begin{table}[H]
\centering
\small
\begin{tabularx}{\linewidth}{Xr}
\toprule
Check & Result\\
\midrule
LOD instruction against a software loop, on the board & 20{,}097 / 20{,}097 values correct\\
Every output equal to the Python integer model (all 192 programs, 960 inputs) & \textbf{960 / 960}\\
Block-RAM cycle counts equal to the Verilator simulation (Speech Commands) & \textbf{240 / 240}\\
Runs finished without bus errors or timeouts & 192 / 192\\
Clock measured by the ARM's timer & 49.997~MHz (50~MHz design)\\
Timing margin at 50~MHz (worst negative slack, WNS) & +1.03~ns (positive: no timing errors)\\
\bottomrule
\end{tabularx}
\caption{Acceptance checks of the board run (\file{fpga\_zybo/results/acceptance.txt}).}
\end{table}

\begin{keyidea}
The real chip computes exactly the same numbers as the Python model and, from block RAM, takes \emph{exactly}
the same number of clock cycles as the simulator, input by input (for example 70{,}271{,}790 cycles for one
spoken word with RoBA-Dn1 int8, in both). So every simulated result in this document holds on real hardware.
\end{keyidea}

\subsection{What the LOD instruction costs in hardware}

Vivado built two versions of the processor: our Ibex with LOD, and the unchanged Ibex without it.

\begin{table}[H]
\centering
\small
\begin{tabularx}{\linewidth}{Xrr}
\toprule
 & with LOD & without LOD\\
\midrule
Ibex core: LUTs & 2{,}471 & 2{,}475\\
Ibex core: flip-flops & 723 & 724\\
Whole design (Ibex + bus + control): LUTs / flip-flops & 4{,}065 / 1{,}916 & --\\
Block RAM & 48.5 of 60 tiles & 48.5 of 60 tiles\\
Timing margin at 50~MHz & +1.03~ns & +1.08~ns\\
Ibex core block (\code{u\_ibex\_core}) power, Vivado estimate & 0.016~W & 0.016~W\\
Whole Ibex (core + register file + bus), Vivado estimate & 0.020~W & 0.019~W\\
\bottomrule
\end{tabularx}
\caption{Area, timing and power from Vivado 2024.2 (\file{fpga\_zybo/reports/}).}
\end{table}

The whole design uses 23\% of the chip's LUTs and 5\% of its flip-flops. The LOD is a handful of gates (the ALU grows, but the synthesis tool reshuffles the rest of the core, so the
total does not change measurably).

\begin{careful}
What these power numbers are, and are not:
\begin{itemize}
  \item 0.016~W is only the Ibex \emph{core} block. The whole Ibex is 0.020~W with LOD and 0.019~W without; the
        192~KB block RAM and the DDR system are counted separately and are not included.
  \item They are Vivado's \emph{estimate} with default activity levels, rounded to 1~mW. They do not depend on
        which program runs, so they \emph{cannot} show whether exact and RoBA draw the same power. ``Same power
        for all methods'' is an assumption, not a measurement.
\end{itemize}
\end{careful}

\subsection{Speed with block RAM and with DDR}

\begin{figure}[H]
\centering
\includegraphics[width=\linewidth]{board_mem}
\caption{Speed-up over exact on the board, Speech Commands. Left: everything in block RAM (same as the
simulator). Right: data in DDR.}
\label{fig:boardmem}
\end{figure}

\begin{figure}[H]
\centering
\includegraphics[width=\linewidth]{board_ddr}
\caption{Speed-up over exact on the board for all three tasks, data in DDR.}
\label{fig:boardddr}
\end{figure}

\begin{table}[H]
\centering
\small
\begin{tabular}{llrrrr}
\toprule
 & & \multicolumn{1}{c}{Speech, BRAM} & \multicolumn{1}{c}{Speech, DDR} & \multicolumn{1}{c}{CIFAR-10, DDR}
 & \multicolumn{1}{c}{ToyADMOS, DDR}\\
\midrule
\multirow{3}{*}{int4} & Square table & 1.10$\times$ & 0.77$\times$ & 0.85$\times$ & 0.86$\times$\\
 & RoBA-Dn1 & 1.20$\times$ & 1.07$\times$ & 1.06$\times$ & 1.04$\times$\\
 & RoBA-Dn2 & 1.06$\times$ & 1.03$\times$ & 1.01$\times$ & 0.99$\times$\\
\midrule
\multirow{3}{*}{int8} & Square table & \textbf{1.95$\times$} & 1.04$\times$ & 1.10$\times$ & 1.25$\times$\\
 & RoBA-Dn1 & 1.81$\times$ & \textbf{1.36$\times$} & \textbf{1.31$\times$} & \textbf{1.40$\times$}\\
 & RoBA-Dn2 & 1.49$\times$ & 1.24$\times$ & 1.19$\times$ & 1.24$\times$\\
\midrule
\multirow{2}{*}{int16} & RoBA-Dn1 & \textbf{2.93$\times$} & \textbf{1.97$\times$} & \textbf{1.95$\times$} & \textbf{2.21$\times$}\\
 & RoBA-Dn2 & 2.39$\times$ & 1.77$\times$ & 1.72$\times$ & 1.89$\times$\\
\bottomrule
\end{tabular}
\caption{Speed-up over exact at the same precision, measured on the board
(\file{fpga\_zybo/results/board\_summary.csv}). Bold: fastest at each precision and memory. The square table
is not feasible at int16 (Section~\ref{sec:tiny}).}
\label{tab:board}
\end{table}

With block RAM the board shows the familiar picture: table first at int8, RoBA far ahead at int16. With DDR,
all speed-ups shrink, and \textbf{the square table falls from first to last}: at int8 RoBA-Dn1 now beats it on
every task, and at int4 the table is even \emph{slower} than exact.

\subsubsection*{Why DDR hurts the table most}

The board counts every DDR access and every cycle the Ibex waits for one. Each access costs about
\textbf{15 cycles} (15.2--15.7 in every run), and waiting took 25--76\% of the running time.

\begin{figure}[H]
\centering
\includegraphics[width=0.9\linewidth]{board_wait}
\caption{Where the time goes with data in DDR (Speech Commands, int8): solid = computing, hatched = waiting for
DDR. In brackets: the same program from block RAM.}
\label{fig:boardwait}
\end{figure}

\begin{example}[: counting the walks to the basement]
For one spoken word at int8, exact, RoBA-Dn1 and RoBA-Dn2 all read the same 5.62~million values from DDR (the
weights and activations), so they all wait the same $\approx$88~million cycles. Their \emph{computing} parts
differ as before (127, 70 and 85 million cycles), so RoBA still saves those cycles, but the speed-up is diluted:
$\dfrac{127 + 88}{70 + 88} = \dfrac{215}{158} = 1.36\times$ instead of $127/70 = 1.81\times$.

The square table computes fastest (65 million) but needs \emph{3.2~million extra} reads, one table look-up per
multiplication, at about 15 cycles each: it waits 141 million cycles instead of 88. Total $65 + 141 = 206$
million, barely better than exact's 214: $1.04\times$.
\end{example}

\begin{keyidea}
A table replaces arithmetic by \emph{memory reads}. That is a great trade when memory is instant, and a bad one
when every read is a walk to the basement. RoBA replaces arithmetic by \emph{a few more instructions on numbers
already in registers}, which costs the same wherever the data lives. On a real microcontroller with slow
external memory, RoBA is the more robust choice.
\end{keyidea}

\subsection{Energy on the board}\label{sec:boardenergy}

Energy is power times time, $E = P \cdot T$. The board measures $T$ exactly (cycles at 50~MHz); $P$ is not
measured, so we give the energy saving under two assumptions and report the pair as a range:

\begin{itemize}
  \item \textbf{equal power}: every method draws the same average power (what Vivado's workload-independent
        estimate suggests). Then the energy saving is simply $1 - 1/\text{speed-up}$;
  \item \textbf{paper's LOD factor}: the core with LOD draws 4.5\% more ($0.829/0.793 = 1.045$), so RoBA's energy
        is multiplied by 1.045. Exact and the square table do not use LOD and keep factor 1.
\end{itemize}

\begin{example}[: energy saving of Dn1 at int16, Speech Commands, block RAM]
The board measured a speed-up of $2.93\times$. Equal power: $1 - 1/2.93 = 65.8\%$ less energy. With the LOD
factor: $1 - 1.045/2.93 = 64.3\%$. So ``64--66\% less energy''.
\end{example}

\begin{table}[H]
\centering
\small
\begin{tabularx}{\linewidth}{lXrrr}
\toprule
 & & & \multicolumn{2}{c}{energy reduction}\\
\cmidrule(lr){4-5}
Precision & Method & board speed-up & equal power & LOD factor 1.045\\
\midrule
int16 & RoBA-Dn1 & 2.93$\times$ & 65.8\% & 64.3\%\\
int16 & RoBA-Dn2 & 2.39$\times$ & 58.2\% & 56.3\%\\
int8 & Square table & 1.95$\times$ & 48.7\% & 48.7\%\\
int8 & RoBA-Dn1 & 1.81$\times$ & 44.9\% & 42.3\%\\
int8 & RoBA-Dn2 & 1.49$\times$ & 32.8\% & 29.7\%\\
int4 & RoBA-Dn1 & 1.20$\times$ & 16.3\% & 12.5\%\\
\bottomrule
\end{tabularx}
\caption{Energy reduction against exact at the same precision, Speech Commands, everything in block RAM, measured
on the board (\file{fpga\_zybo/results/board\_summary.csv}, columns \code{energy\_vs\_exact\_fpga} and
\code{energy\_vs\_exact}).}
\label{tab:boardenergy}
\end{table}

With the data in DDR the reductions are smaller, because the waiting for memory is the same for every method:

\begin{table}[H]
\centering
\small
\begin{tabular}{lrrr}
\toprule
int16 RoBA-Dn1, data in DDR & board speed-up & equal power & LOD factor 1.045\\
\midrule
CIFAR-10 & 1.95$\times$ & 48.8\% & 46.4\%\\
Speech Commands & 1.97$\times$ & 49.2\% & 46.8\%\\
ToyADMOS & 2.21$\times$ & 54.7\% & 52.6\%\\
\bottomrule
\end{tabular}
\caption{Energy reduction of RoBA-Dn1 at int16 with the data in DDR, measured on the board.}
\end{table}

\begin{careful}
Do not mix these board numbers with the simulation numbers of Section~\ref{sec:tiny} in one table. For example,
the $3.04\times$ / $1.77\times$ / $1.41\times$ of Section~\ref{sec:tiny} are CIFAR-10 in the Verilator study
(int16 Dn1, int8 Dn1, int8 Dn2); the board's block-RAM numbers above are Speech Commands. Both are correct, but
for different tasks and builds, so label each as simulation or board.
\end{careful}

\begin{keyidea}
In one sentence, as it could go in a paper: \emph{Estimated energy reduction, $E = P\cdot T$, assuming the same
average power for all methods (Vivado's estimate: 16~mW for the Ibex core with or without LOD,
workload-independent). With the paper's LOD power factor of 1.045 instead, the reductions are 1.5--3 points lower
at int8 and int16; for example 64--66\% for RoBA-Dn1 at int16.} The conclusion does not depend on which power
assumption is used.
\end{keyidea}

Replacing the equal-power assumption by a power figure per method would need Vivado's power analysis with the
real switching activity of each program (from a simulation of the design running it); that is listed under
future work.

\subsection{What Phase 4 adds to the story}

\begin{enumerate}
  \item \textbf{The simulation is trustworthy.} The real chip gives identical outputs (960/960) and, from block
        RAM, identical cycle counts (240/240).
  \item \textbf{The LOD instruction is practically free in hardware}: no measurable area, the core's estimated
        power unchanged at 1~mW resolution, and the design meets timing at 50~MHz with margin.
  \item \textbf{Energy}: on the board RoBA-Dn1 at int16 needs 64--66\% less energy than exact from block RAM
        and 46--55\% less with data in DDR (range = the two power assumptions).
  \item \textbf{Memory speed changes the ranking.} With fast on-chip memory the table wins at int8; with
        external DDR, RoBA-Dn1 wins at int8 (1.31--1.40$\times$) and at int16 (1.95--2.21$\times$), and the
        table drops to about exact speed or worse.
\end{enumerate}
